% Figures de l'étude 02 (quantification extrême des poids), en pgfplots.
% Chaque série est un fichier fig-*.csv commité, produit par outils/quant-csv-memoire.py
% à partir des journaux ; aucune valeur n'est écrite à la main ici.
% Le document hôte charge pgfplots (compat >= 1.18) et peut redéfinir \CheminQuant
% (dossier csv/, relatif au dossier de compilation).
% Sous lualatex, les barres d'erreur exigent « lua backend=false ».
\providecommand{\CheminQuant}{../fondations/journaux/02-quantification/csv}
\pgfplotstableset{col sep=comma}
\pgfplotsset{quant/.style={lua backend=false, grid=major, legend cell align=left,
  every axis plot/.append style={line width=0.8pt}},
  quant erreurs/.style={error bars/.cd, y dir=both, y explicit}}

% Figure Q1 : distorsion contre débit sur la source N(0,1), mesuré ici (§ 2).
\newcommand{\FigureQuantMSE}{%
\begin{tikzpicture}
\begin{semilogyaxis}[quant, width=0.92\linewidth, height=0.62\linewidth,
  xlabel={bits par poids $R$}, ylabel={MSE pour $\sigma^2 = 1$},
  xmin=0.9, xmax=3.1, ymin=0.012, ymax=0.5, legend pos=north east,
  log ticks with fixed point, ytick={0.02,0.05,0.1,0.2,0.4}]
\addplot[black, no markers, domain=0.9:3.1, samples=50] {2^(-2*x)};
\addlegendentry{borne de Shannon $2^{-2R}$}
\addplot[only marks, mark=square*, quant erreurs] table[x=bits, y=mse, y error=et, col sep=comma] {\CheminQuant/fig-q1-lloyd.csv};
\addlegendentry{Lloyd-Max (scalaire optimal)}
\addplot[only marks, mark=triangle*, mark size=2.5pt, quant erreurs] table[x=bits, y=mse, y error=et, col sep=comma] {\CheminQuant/fig-q1-vq8.csv};
\addlegendentry{VQ dim.~8, livre aléatoire}
\addplot[only marks, mark=diamond*, mark size=2.8pt, quant erreurs] table[x=bits, y=mse, y error=et, col sep=comma] {\CheminQuant/fig-q1-iq2xxs.csv};
\addlegendentry{IQ2\_XXS (llama.cpp)}
\addplot[only marks, mark=*, quant erreurs] table[x=bits, y=mse, y error=et, col sep=comma] {\CheminQuant/fig-q1-treillis.csv};
\addlegendentry{treillis à décalage, $L=16$}
\end{semilogyaxis}
\end{tikzpicture}}

% Figure Q2 : produit matrice-vecteur 16 384 x 16 384 sur la RTX 4060, Go/s équivalents
% (médianes ; Vulkan passe 4, DirectX 12 passe 3). Plafond VRAM 256 Go/s.
\newcommand{\FigureQuantGEMV}{%
\begin{tikzpicture}
\begin{axis}[quant, width=0.92\linewidth, height=0.55\linewidth, ybar, bar width=9pt,
  ylabel={Go/s équivalents}, ymin=0, ymax=290, xmin=-0.6, xmax=5.6,
  xtick={0,...,5}, xticklabels from table={\CheminQuant/fig-q2-vulkan.csv}{format},
  x tick label style={rotate=25, anchor=east, font=\small}, legend pos=north east,
  nodes near coords, nodes near coords style={font=\tiny}, major x grid style={draw=none}]
\addplot table[x=i, y=go_s, col sep=comma] {\CheminQuant/fig-q2-vulkan.csv};
\addlegendentry{Vulkan}
\addplot table[x=i, y=go_s, col sep=comma] {\CheminQuant/fig-q2-dx12.csv};
\addlegendentry{DirectX 12}
\draw[dashed] (axis cs:-0.6,256) -- (axis cs:5.6,256);
\end{axis}
\end{tikzpicture}}

% Figure Q3 : perplexité contre octets lus par jeton, Qwen3.5 4B et 9B, mesuré ici (§ 5).
\newcommand{\FigureQuantEchelle}{%
\begin{tikzpicture}
\begin{axis}[quant, width=0.92\linewidth, height=0.62\linewidth,
  xlabel={Go lus par jeton}, ylabel={perplexité (wikitext-2, 40 $\times$ 512 jetons)},
  xmin=1.2, xmax=5.6, ymin=6.5, ymax=16.5, legend pos=north east,
  nodes near coords, point meta=explicit symbolic,
  nodes near coords style={font=\scriptsize, anchor=south west}]
\addplot[only marks, mark=*, quant erreurs] table[x=go, y=ppl, y error=et, meta=etiquette, col sep=comma] {\CheminQuant/fig-q3-4B.csv};
\addlegendentry{Qwen3.5-4B}
\addplot[only marks, mark=square*, nodes near coords style={anchor=north west}, quant erreurs] table[x=go, y=ppl, y error=et, meta=etiquette, col sep=comma] {\CheminQuant/fig-q3-9B.csv};
\addlegendentry{Qwen3.5-9B}
\end{axis}
\end{tikzpicture}}

% Figure Q4 : à octets égaux, Llama 2 en treillis QTIP (perplexité ; littérature, QTIP tab. 5).
\newcommand{\FigureQuantIsoOctets}{%
\begin{tikzpicture}
\begin{semilogxaxis}[quant, width=0.92\linewidth, height=0.6\linewidth,
  xlabel={Go lus par jeton (échelle log.)}, ylabel={perplexité WikiText-2, contexte 4096},
  log ticks with fixed point, xtick={2,4,8,16,32}, xmin=1.4, xmax=40, legend pos=north east, nodes near coords, point meta=explicit symbolic,
  nodes near coords style={font=\scriptsize, anchor=west}]
\addplot[only marks, mark=*] table[x=go, y=ppl, meta=etiquette, col sep=comma] {\CheminQuant/fig-q4-llama2-7B.csv};
\addlegendentry{Llama 2 7B}
\addplot[only marks, mark=square*] table[x=go, y=ppl, meta=etiquette, col sep=comma] {\CheminQuant/fig-q4-llama2-13B.csv};
\addlegendentry{Llama 2 13B}
\addplot[only marks, mark=triangle*, mark size=2.5pt] table[x=go, y=ppl, meta=etiquette, col sep=comma] {\CheminQuant/fig-q4-llama2-70B.csv};
\addlegendentry{Llama 2 70B}
\end{semilogxaxis}
\end{tikzpicture}}

% Figure Q6 : vitesse de décodage (lot de 1, tout en VRAM, CUDA) contre octets lus par jeton,
% mesuré ici (passe 3) ; la droite pointillée est le plafond 256 Go/s.
\newcommand{\FigureQuantVitesse}{%
\begin{tikzpicture}
\begin{axis}[quant, width=0.92\linewidth, height=0.62\linewidth,
  xlabel={Go lus par jeton}, ylabel={jetons par seconde},
  xmin=1.2, xmax=5.6, ymin=0, ymax=180, legend pos=north east,
  nodes near coords, point meta=explicit symbolic,
  nodes near coords style={font=\scriptsize, anchor=south west}]
\addplot[black, dashed, no markers, domain=1.2:5.6, samples=60] {256/x};
\addlegendentry{plafond $256/B$}
\addplot[mark=*] table[x=go, y=jetons_s, meta=etiquette, col sep=comma] {\CheminQuant/fig-q6-vitesse-4B.csv};
\addlegendentry{Qwen3.5-4B}
\addplot[mark=square*] table[x=go, y=jetons_s, meta=etiquette, col sep=comma] {\CheminQuant/fig-q6-vitesse-9B.csv};
\addlegendentry{Qwen3.5-9B}
\end{axis}
\end{tikzpicture}}

% Figure Q5 : à octets égaux, Qwen3 au MMLU (littérature, Tetra tab. 3 ; Go = tampons sur la carte).
\newcommand{\FigureQuantQwenMMLU}{%
\begin{tikzpicture}
\begin{semilogxaxis}[quant, width=0.92\linewidth, height=0.6\linewidth,
  xlabel={Go de poids lus par jeton (échelle log.)}, ylabel={MMLU (14 042 questions)},
  log ticks with fixed point, xtick={1,2,4,8,16,32}, xmin=1, xmax=40, legend pos=south east, nodes near coords, point meta=explicit symbolic,
  nodes near coords style={font=\scriptsize, anchor=south}]
\addplot[mark=*] table[x=go, y=mmlu, meta=etiquette, col sep=comma] {\CheminQuant/fig-q4-qwen3-mmlu-tetra.csv};
\addlegendentry{Tetra, $\approx 2{,}7$ bits/paramètre}
\addplot[mark=square*] table[x=go, y=mmlu, meta=etiquette, col sep=comma] {\CheminQuant/fig-q4-qwen3-mmlu-awq.csv};
\addlegendentry{AWQ 4 bits, 5,3 à 6,0 bits/paramètre}
\addplot[mark=triangle*, mark size=2.5pt] table[x=go, y=mmlu, meta=etiquette, col sep=comma] {\CheminQuant/fig-q4-qwen3-mmlu-fp16.csv};
\addlegendentry{FP16}
\end{semilogxaxis}
\end{tikzpicture}}
